% generated table
% source run: ig_examples_200pos_v1/attribution_metrics_draws20
% Selection and wording follow the part B memo behind issue #120; the numbers come from the
% run stamped above (issue #187 re-sample). Regenerate from that summary rather than
% editing the numbers here.
\begin{table}[t]
\centering
\small
\setlength{\tabcolsep}{4pt}
\begin{tabular}{llrrrr}
\toprule
& & \multicolumn{2}{c}{$\rho_{\text{LOO}}$} & \multicolumn{2}{c}{DelAUC gap} \\
\cmidrule(lr){3-4}\cmidrule(lr){5-6}
Model & Method & base & ABTT & base & ABTT \\
\midrule
LaTa & IG & 0.003 & \textbf{0.370} & \textbf{0.859} & 0.181 \\
 & MaRC & 0.083 & \textbf{0.538} & \textbf{0.524} & 0.327 \\
\addlinespace[2pt]
PhilTa & IG & 0.329 & \textbf{0.614} & 0.113 & \textbf{0.403} \\
 & MaRC & \textbf{0.469} & 0.445 & 0.201 & \textbf{0.313} \\
\addlinespace[2pt]
mT5-base & IG & 0.143 & \textbf{0.686} & 0.061 & \textbf{0.561} \\
 & MaRC & 0.257 & \textbf{0.504} & 0.042 & \textbf{0.463} \\
\bottomrule
\end{tabular}
\caption{Attribution faithfulness at the predeclared operational layers, 200 positive pairs per model, for integrated gradients (IG) and retrieval-adapted MaRC. $\rho_{\text{LOO}}$ correlates attribution magnitude with the leave-one-out change in the cosine. DelAUC gap is the random-order minus attribution-order deletion-curve area, positive when attribution beats chance; the reference runs from 0.692 to 0.961 here, so zero is chance. Higher is better in both; boldface marks the better variant. ABTT wins 5/6 and 4/6. Ties, within 2 standard errors of zero: $\rho$ for PhilTa MaRC. Ratio metrics are undefined below a full-query cosine of 0.05, so the DelAUC columns average 195 to 199 ABTT pairs against 199 to 200 baseline pairs. Cross-variant comparisons are descriptive. Secondary metrics: Table~\ref{tab:attribution_metrics_secondary}.}
\label{tab:attribution_metrics_main}
\end{table}
